\section[Photodetector Architectures]{Photodetector Architectures}
\label{sec:pd_architectures}

With the metrics of \cref{sec:pd_metrics} as the evaluation framework, the
landscape of photodetector device classes can be surveyed systematically. Each
architecture is assessed against responsivity, bandwidth, and noise, and each
is found to resolve one specific limitation of its predecessor, until the
waveguide-integrated PIN junction emerges not as a historical preference but as
the consequence of removing every identified obstacle in sequence.

\subsection[PN and PIN Junctions]{PN and PIN Junctions: Depletion Engineering}
\label{subsec:pd_pn_pin}

\paragraph{Photoconductive Detection.}
The photoconductor is the simplest optical transducer because it requires no
junction and no depletion engineering. Its apparent advantage is internal gain:
\begin{equation}
    G_{\mathrm{PC}} = \frac{\tau_r}{\tau_{\mathrm{tr}}},
    \qquad
    \tau_{\mathrm{tr}} = \frac{L^2}{\mu_n V_b},
    \label{eq:pc_gain}
\end{equation}
where $\tau_r$ is the minority-carrier recombination lifetime, $\tau_{\mathrm{tr}}$
is the carrier transit time across the contact separation $L$, $\mu_n$ is the
electron mobility, and $V_b$ is the applied bias. Physically, $G_{\mathrm{PC}}$
counts how many complete circuits a majority carrier traverses before
recombination removes its partner: if a photogenerated hole survives for
many transit times, the external circuit must replenish the missing electron
repeatedly, yielding an apparent multiplication of the photocurrent. This gain
mechanism, however, is inseparable from the recombination lifetime that sustains
it. Because the same $\tau_r$ that produces gain also governs the impulse
response, the $-3\,\mathrm{dB}$ bandwidth is:
\begin{equation}
    f_{3\,\mathrm{dB,PC}} = \frac{1}{2\pi\tau_r}.
    \label{eq:pc_bw}
\end{equation}
The gain-bandwidth product $G_{\mathrm{PC}}\times f_{3\,\mathrm{dB,PC}}$ is
therefore a constant set by material parameters alone, with no free geometric
lever. For germanium, trap-limited lifetimes in the range
$\tau_r\sim\SI{1}{\micro\second}$ constrain the bandwidth to the
sub-megahertz regime, placing the photoconductor many orders of magnitude below
the \SI{50}{\giga\hertz} electrical bandwidth required for
\SI{100}{\giga\baud}-class operation in the target data-centre interconnect
context. The photoconductor is therefore rejected by this intrinsic constraint
before any geometric optimisation is considered.

\paragraph{PN Junction Photodiodes.}
The PN junction eliminates the most critical weakness of the photoconductor
by replacing the ohmic resistor with a rectifying interface that sustains
an internal built-in electric field $E_{\mathrm{bi}}$ across the depletion
region. Photogenerated electron-hole pairs created within, or diffusing into,
this depletion layer experience a drift force that sweeps electrons toward
the n-contact and holes toward the p-contact at velocities approaching the
material saturation velocity $v_{\mathrm{sat}}$. The recombination lifetime
no longer limits speed because carriers are extracted before recombination can
occur, decoupling the bandwidth from $\tau_r$.

The drift component of the frequency response is well approximated, for a
uniformly illuminated depletion region of width $W_d$, by:
\begin{equation}
    H_{\mathrm{drift}}(f) =
    \frac{1 - e^{-j2\pi f\,W_d/v_{\mathrm{sat}}}}{j2\pi f\,W_d/v_{\mathrm{sat}}},
    \label{eq:hdrift}
\end{equation}
which is the Fourier representation of a uniform current pulse of duration
$\tau_{\mathrm{tt}} = W_d/v_{\mathrm{sat}}$. The $-3\,\mathrm{dB}$ point of
$|H_{\mathrm{drift}}|$ occurs near $f_{\mathrm{tt}} \approx 0.45\,v_{\mathrm{sat}}/W_d$,
consistent with \cref{eq:bw_tt}.

A residual limitation arises because not all photogeneration is confined to the
depletion layer. Carriers created in the quasi-neutral p- and n-regions must
diffuse to the depletion edge before the field can extract them. This diffusion
process introduces a slow tail in the impulse response, characterised by the
ambipolar diffusion time constant $\tau_D$. The composite transfer function
acquires a low-pass factor:
\begin{equation}
    H_{\mathrm{PN}}(f) =
    H_{\mathrm{drift}}(f)\,
    \frac{1}{1+j2\pi f\tau_D},
    \label{eq:hpn}
\end{equation}
where $\tau_D$ is determined by the minority-carrier diffusion length and
the distance from the photogeneration centroid to the depletion edge. When
the diffusion-limited pole $1/(2\pi\tau_D)$ falls below the drift-limited
transit-time bandwidth, it becomes the binding constraint on the overall
detector speed. In practice, this condition is met at moderate symbol rates,
but at the \SIrange{50}{100}{\giga\hertz} electrical bandwidth class targeted
here, even a modest quasi-neutral absorption fraction is sufficient to suppress
the frequency response to an unacceptable degree. The PN structure therefore
resolves the gain-bandwidth lock of the photoconductor while retaining a
diffusion-tail bottleneck that must be addressed by the next architectural step.

\paragraph{PIN Photodiodes.}
The PIN structure eliminates the diffusion tail by inserting a lightly doped,
or intrinsic, semiconductor layer of thickness $W_i$ between the p- and n-doped
contact regions. Under reverse bias $V_{\mathrm{bias}}$, the low doping density
of the intrinsic layer ensures that it is fully depleted at modest voltages,
because the depletion depth in a one-sided junction scales as
$W \propto \sqrt{(V_{\mathrm{bi}}+|V_{\mathrm{bias}}|)/N_a}$, so a very low
effective doping $N_a$ in the intrinsic region allows full depletion to be
achieved without requiring large voltages. Once fully depleted, the electric
field across the intrinsic region is approximately uniform:
\begin{equation}
    E(x) \approx \frac{|V_{\mathrm{bias}}|+V_{\mathrm{bi}}}{W_i},
    \qquad x\in[0,W_i],
    \label{eq:pin_efield_uniform}
\end{equation}
where $V_{\mathrm{bi}}$ is the built-in potential determined by the doping
contrast at the p--i and i--n interfaces. This expression follows directly from
integrating Poisson's equation across a region of negligible space charge:
$\mathrm{d}E/\mathrm{d}x \approx 0$, so the field is linear only at the
boundaries and constant throughout the bulk of the intrinsic layer. The
practical importance of this uniformity is that all optically generated
carriers within $W_i$ experience essentially the same extraction field,
eliminating the position-dependent transit-time dispersion that would
arise in a non-uniform field profile.

If the absorbing material is selected so that its absorption coefficient
$\alpha$ at the operating wavelength satisfies $\alpha^{-1}\lesssim W_i$,
the overwhelming majority of photogeneration occurs inside the fully depleted
intrinsic region. Diffusion from quasi-neutral zones is thereby suppressed,
and the diffusion tail of \cref{eq:hpn} is eliminated. The bandwidth is
then governed entirely by the two competing mechanisms already formalised
through \cref{eq:bw_tt,eq:bw_rc}:
\begin{equation}
    \frac{1}{f_{3\,\mathrm{dB}}^2} \approx
    \frac{1}{f_{\mathrm{tt}}^2} + \frac{1}{f_{\mathrm{RC}}^2},
    \label{eq:pin_bw_combined}
\end{equation}
where $f_{\mathrm{tt}} \propto v_{\mathrm{sat}}/W_i$ and
$f_{\mathrm{RC}} \propto (R_s C_j)^{-1}$ with $C_j \propto A/W_i$. An
increase in $W_i$ raises the transit time $\tau_{\mathrm{tt}}$ and thus
degrades $f_{\mathrm{tt}}$, while simultaneously reducing the junction
capacitance $C_j$ and thus improving $f_{\mathrm{RC}}$. There is therefore
an optimum intrinsic thickness $W_i^{\star}$ that maximises the total
bandwidth for a given junction area $A$ and series resistance $R_s$:
\begin{equation}
    W_i^{\star} = \left(\frac{v_{\mathrm{sat}}\,A\,\varepsilon_0\varepsilon_r}
    {2\pi\,R_s}\right)^{1/2},
    \label{eq:pin_opt_wi}
\end{equation}
obtained by differentiating \cref{eq:pin_bw_combined} with respect to $W_i$
and equating the transit-time and RC contributions at the optimum. The
structural advantage of the PIN over the PN diode is that $W_i$ is a
lithographically controlled epitaxial parameter, decoupled from the contact
doping levels used to minimise $R_s$. This additional degree of freedom---the
independent control of the depletion thickness---is the key reason the PIN
remains the canonical high-speed receiver device: it permits the junction
geometry to be engineered at a level that the simple PN diode, where $W_d$
is set by the doping profile alone, cannot match.

\subsection[Avalanche Photodiodes]{Avalanche Photodiodes: Gain versus Noise}
\label{subsec:pd_apd}

An avalanche photodiode (APD) exploits impact ionisation within a high-field
multiplication region to amplify the photocurrent internally, yielding an
effective responsivity:
\begin{equation}
    \mathcal{R}_{\mathrm{APD}} = M\,\mathcal{R}_{\mathrm{PIN}},
    \label{eq:r_apd}
\end{equation}
where $M$ is the mean avalanche multiplication factor and
$\mathcal{R}_{\mathrm{PIN}}$ is the responsivity that the same absorption
layer would provide without multiplication. At first inspection this appears
to offer a route to higher sensitivity without increasing the optical input
power, but the avalanche mechanism introduces two penalties that are absent
in the low-bias PIN.

The first is a fundamental gain-bandwidth constraint. Avalanche multiplication
is not instantaneous: the carriers initiating an impact-ionisation event must
traverse the high-field region, and the secondary carriers they generate
must in turn complete their own transit before the multiplication chain
terminates. The statistical build-up and decay of this cascade introduces
a finite temporal spread that broadens the impulse response. For a multiplication
region of length $l_{\mathrm{mult}}$ and saturated ionisation velocity
$v_{\mathrm{ion}}$, the effective multiplication bandwidth is approximately:
\begin{equation}
    f_M \approx \frac{v_{\mathrm{ion}}}{2\pi\,k\,M\,l_{\mathrm{mult}}},
    \label{eq:apd_gbp}
\end{equation}
where $k = \alpha_h/\alpha_e$ is the ratio of hole to electron impact-ionisation
coefficients. The product $M\times f_M$ is therefore a material- and
geometry-specific constant, conventionally termed the gain-bandwidth product
(GBP). As $M$ is increased to boost sensitivity, $f_M$ decreases
proportionally, and the bandwidth headroom required for the target
\SI{50}{\giga\hertz} class rapidly exhausts the available GBP. In germanium,
the nearly equal ionisation coefficients ($k\approx 1$) make this constraint
particularly severe compared with materials such as InAlAs, where $k\ll 1$
permits higher GBP.

The second penalty is excess noise. Even in the absence of thermal noise,
the stochastic nature of the multiplication process introduces fluctuations
in $M$ about its mean value that manifest as an additional noise current.
The McIntyre model~\cite{McIntyre1966} characterises this through an
excess noise factor $F(M)$, which modifies the shot-noise power spectral
density from the unmultiplied case:
\begin{equation}
    S_{I,\mathrm{APD}} = 2q\,I_{\mathrm{ph}}\,M^2\,F(M)\,\Delta f,
    \label{eq:apd_noise}
\end{equation}
where $I_{\mathrm{ph}}$ is the primary photocurrent and $\Delta f$ is the
noise bandwidth. For a single-carrier-initiated avalanche with ionisation
ratio $k$, McIntyre showed that:
\begin{equation}
    F(M) = k\,M + (1-k)\!\left(2 - \frac{1}{M}\right).
    \label{eq:apd_fn}
\end{equation}
In the ideal limit $k\rightarrow 0$ (electrons ionise, holes do not),
$F\rightarrow 2$ as $M\rightarrow\infty$, and the noise penalty is modest.
In the opposite regime $k\rightarrow 1$, which approximates the situation
in Ge where $\alpha_e\approx\alpha_h$, \cref{eq:apd_fn} reduces to
$F\approx M$, and substitution into \cref{eq:apd_noise} shows that the
noise power spectral density grows as $M^3$ rather than $M^2$. The
net signal-to-noise benefit of multiplication then degrades rapidly with
increasing $M$, and the optimum gain that minimises the noise-equivalent
power saturates at values well below those that would be useful for
sensitivity improvement in a thermal-noise-dominated receiver front-end.

Taking both penalties together---the GBP constraint and the excess noise
factor---the APD in Ge offers no practical advantage over a low-bias PIN
receiver paired with a low-noise transimpedance amplifier (TIA) at the
data rates of interest. The TIA thermal noise, which sets the dominant
noise floor in high-speed receivers, can be reduced by circuit design
without introducing the bandwidth restriction inherent to avalanche
multiplication. The APD is therefore not selected for the present work,
and the PIN operating at modest reverse bias is retained as the active
device architecture.

\subsection[Waveguide-Integrated Photodetectors]{Waveguide-Integrated Photodetectors: Interaction Length Advantage}
\label{subsec:pd_waveguide}

The PIN junction resolves the device-level bandwidth problem but does not, by
itself, resolve the geometric conflict between absorption efficiency and transit
time. This conflict is most acute in a surface-illuminated (or
normal-incidence) configuration, where the optical beam enters through the top
surface and propagates vertically through the absorbing layer. In this geometry
the same dimension, the absorber thickness $W_i$, must simultaneously satisfy
two contradictory requirements: the Beer-Lambert absorption condition,
$W_i \gtrsim \alpha^{-1}$, to achieve high quantum efficiency, and the transit-time
bandwidth condition, $W_i \ll v_{\mathrm{sat}}/(2\pi f_{3\,\mathrm{dB}})$,
derived from \cref{eq:bw_tt}. At \SI{1310}{\nano\metre}, the auxiliary MATLAB
model uses $\alpha\approx\SI{7e3}{\per\centi\metre}$ for germanium, corresponding
to a bulk $1/e$ absorption depth of $\alpha^{-1}\approx\SI{1.4}{\micro\metre}$
before surface-normal reflection and confinement penalties are included.
Achieving 90\,\% absorption requires a several-micrometre optical path, while
achieving a transit-time-limited bandwidth of \SI{50}{\giga\hertz} requires
$W_i \lesssim \SI{1}{\micro\metre}$ for typical saturation velocities in Ge
($v_{\mathrm{sat}}\approx\SI{6e6}{\centi\metre\per\second}$). These two
requirements differ by over an order of magnitude, and no single value of
$W_i$ can satisfy both simultaneously. The surface-illuminated PIN geometry
is therefore fundamentally overconstrained for high-speed, high-efficiency
operation in the O-band.

Waveguide integration resolves this conflict by a geometric transformation:
the optical mode is laterally confined by the waveguide structure and
propagates longitudinally along the length of the detector rather than
vertically through its thickness. The relevant dimension for optical absorption
is then the propagation length $L$ along the device axis, while the vertical
dimension $W_i$ retains its role as the carrier collection thickness that
determines the transit time. Because $L$ and $W_i$ are orthogonal spatial
coordinates, they can be engineered independently.

The evolution of optical power along the waveguide-integrated absorber follows
the Beer-Lambert law in the propagation direction:
\begin{equation}
    P(z) = P_0\,e^{-\Gamma\alpha z},
    \label{eq:wg_beer_lambert}
\end{equation}
where $P_0$ is the power entering the absorber at $z=0$, $\alpha$ is the
material absorption coefficient of germanium, and $\Gamma$ is the optical
confinement factor, defined as the fraction of the guided modal power that
overlaps spatially with the germanium absorbing region:
\begin{equation}
    \Gamma = \frac{\displaystyle\int_{\mathrm{Ge}}
        |\mathbf{E}(\mathbf{r}_\perp)|^2\,\mathrm{d}A_\perp}
    {\displaystyle\int_{-\infty}^{\infty}
        |\mathbf{E}(\mathbf{r}_\perp)|^2\,\mathrm{d}A_\perp},
    \label{eq:mode_overlap}
\end{equation}
with the integration in the numerator restricted to the transverse cross-section
of the Ge layer and the denominator taken over the full modal field. The
factor $\Gamma$ is determined by the waveguide geometry: for a Ge-on-Si
butt-coupled or evanescent-coupling scheme, $\Gamma$ can range from values
approaching unity in a strongly confining geometry to values well below
0.5 in an evanescent-coupling taper architecture where the mode is
gradually transferred from the Si waveguide into the Ge absorber.

The total absorption efficiency, relating absorbed power to incident power,
is obtained by integrating \cref{eq:wg_beer_lambert} over the device length $L$:
\begin{equation}
    \eta_{\mathrm{abs}} = \frac{P_0 - P(L)}{P_0} = 1 - e^{-\Gamma\alpha L}.
    \label{eq:wg_abs_efficiency}
\end{equation}
This expression reveals the decisive advantage of the waveguide architecture:
$\eta_{\mathrm{abs}}$ can be driven arbitrarily close to unity by increasing
$L$, entirely independently of $W_i$. For a representative value of
$\Gamma = 0.8$ and $\alpha = \SI{7e3}{\per\centi\metre}$ at
\SI{1310}{\nano\metre}, a several-micrometre waveguide absorber can achieve
high absorption, while $W_i$ remains freely
settable in the sub-micron regime required for high transit-time bandwidth.
The geometric orthogonality between $L$ and $W_i$ is therefore complete:
$L$ governs responsivity through \cref{eq:wg_abs_efficiency}, while $W_i$
governs $f_{\mathrm{tt}}$ through \cref{eq:bw_tt}, and the two parameters
are adjusted in a two-dimensional design space without mutual constraint.

A secondary benefit of the waveguide geometry is the suppression of the
RC penalty associated with a large device footprint. In the surface-illuminated
case, a large area is needed to intercept the incident beam, increasing $C_j$
and degrading $f_{\mathrm{RC}}$. In the waveguide-integrated device, the
transverse cross-section of the mode is set by the single-mode waveguide
dimensions, typically of order $\SI{0.5}{\micro\metre}\times\SI{0.5}{\micro\metre}$
in sub-micron Si photonic platforms. The junction width $W$ and height $W_i$
are therefore bounded by the mode confinement rather than by beam-capture
requirements, enabling small junction areas and correspondingly small
capacitances. The longitudinal length $L$ does increase the total capacitance
as $C_j\propto W\cdot L/W_i$, but because $W$ is sub-micron, the product
$W\cdot L$ for a \SI{10}{\micro\metre}--\SI{20}{\micro\metre}-long device
remains well below the junction areas necessary in surface-illuminated
geometries for equivalent coupling efficiency.

\begin{figure}[H]
    \centering
    \missingfigure[figwidth=0.88\linewidth]{
        Comparative schematic of surface-normal illumination and waveguide
        coupling, highlighting the distinction between optical propagation
        length $L$ and depletion thickness $W_i$ and the design freedom gained
        when these two dimensions are separated.
    }
    \caption{Waveguide coupling is the decisive architectural refinement because
    it decouples the absorption dimension from the transit-time dimension. That
    geometric change turns an overconstrained surface-normal problem into a
    two-parameter design space.}
    \label{fig:wg_vs_sni}
\end{figure}

\noindent The geometric orthogonality introduced by waveguide integration
does not eliminate all design tension: as $L$ is increased to improve
$\eta_{\mathrm{abs}}$, the total junction capacitance grows proportionally,
degrading $f_{\mathrm{RC}}$. The optimum device length is therefore set not
by \cref{eq:wg_abs_efficiency} in isolation but by the joint optimisation of
the bandwidth expression in \cref{eq:pin_bw_combined}, in which $C_j$ now
depends on $L$ as well as on $W_i$ and $W$. This residual coupling between
absorption efficiency and RC bandwidth---absent from the transit-time
discussion but present in the full bandwidth budget---is the central
design tension addressed in \cref{sec:pd_design}.

\subsection[Selection Rationale]{Selection Rationale: Why PIN for High-Speed Integrated Links}
\label{subsec:pd_arch_compare}

The four preceding device classes can be evaluated against a unified set of
performance criteria---responsivity, bandwidth, and noise---by tracing the
physical mechanism that limits each architecture and identifying the structural
modification that removes it. \Cref{fig:pd_evolution_ladder} condenses this
selection logic from the photoconductor through to the final architecture.

\begin{figure}[H]
    \centering
    \missingfigure[figwidth=0.88\linewidth]{
        Comparative schematic linking the photoconductor, PN photodiode, PIN
        photodiode, and APD to the governing physical arguments that motivate
        the final waveguide-integrated Ge-on-Si PIN architecture.
    }
    \caption{Architectural justification of the waveguide-integrated Ge-on-Si
    PIN photodetector. Each comparison resolves a specific limitation of the
    preceding alternative, leaving the waveguide-integrated Ge-on-Si PIN as the
    architecture compatible with the O-band, low-voltage, high-bandwidth target
    of this chapter.}
    \label{fig:pd_evolution_ladder}
\end{figure}

The photoconductor is rejected because the gain-bandwidth product
$G_{\mathrm{PC}}\times f_{3\,\mathrm{dB,PC}}$, fixed by the recombination
lifetime via \cref{eq:pc_bw}, cannot reach the target bandwidth without
eliminating the gain entirely. The PN junction removes the recombination
constraint but retains a diffusion-tail bandwidth penalty from quasi-neutral
absorption, as expressed by the pole in \cref{eq:hpn}. The PIN structure
eliminates the diffusion tail by confining all useful photogeneration within
the fully depleted intrinsic layer, reducing the bandwidth competition to
the transit-time and RC terms of \cref{eq:pin_bw_combined}, both of which
are geometrically controllable. The APD appears to restore sensitivity
through multiplication but introduces a gain-bandwidth product constraint
via \cref{eq:apd_gbp} and an excess noise penalty via \cref{eq:apd_fn}
that, for germanium with $k\approx 1$, grows as $F\approx M$ and rapidly
erases the SNR benefit; the bandwidth headroom needed for the target data
rate is not available at any useful gain.

The waveguide-integrated PIN is therefore not simply the preferred architecture
but the inevitable consequence of eliminating each identified physical
limitation in sequence. It provides independent control of the two parameters
governing responsivity and bandwidth, namely $L$ and $W_i$ respectively,
through a geometric orthogonality that is unavailable in any of the preceding
device classes. Operating at low reverse bias, it avoids the excess noise
and reliability concerns of avalanche operation, and it presents a junction
capacitance commensurate with the narrow transverse dimensions of the
Si photonic waveguide mode. The dominant noise source in the resulting
receiver is the thermal noise of the external transimpedance amplifier,
which is addressed at the circuit level without compromising the
photodetector bandwidth. The remaining design problem---achieving simultaneously
high responsivity and high bandwidth within this waveguide-integrated PIN
framework at \SI{1310}{\nano\metre} in Ge-on-Si---is the tension addressed
by the material analysis of \cref{sec:pd_material}, the electrode and geometry
engineering of \cref{sec:pd_design}, and the literature survey of
\cref{sec:pd_literature}.
